\documentclass[letterpaper,11pt]{article}
\usepackage[top=0.8in, bottom=0.8in, left=0.9in, right=0.9in]{geometry}
\usepackage[hidelinks]{hyperref}
\usepackage{lmodern}
\usepackage{parskip}
\usepackage{microtype}
\setlength{\parskip}{6pt}
\setlength{\parindent}{0pt}
\begin{document}
\begin{flushleft}
\textbf{\large Ajay Venkatesh} \\
Brooklyn, NY \\
+1 516 585 9013 \\
\href{mailto:av3855@nyu.edu}{av3855@nyu.edu}
\end{flushleft}
\vspace{10pt}
May 21, 2026
\vspace{6pt}

Hiring Manager \\
University of Michigan

\vspace{10pt}

Dear Hiring Manager,

My name is Ajay Venkatesh. I'm finishing my M.S. in Computer Engineering at NYU, with production software experience at PwC in Bengaluru, a summer SDE internship at Amazon, and ongoing on-campus work at NYU's Novel AI Technologies. I'm writing about the Data Scientist role at the University of Michigan.

The role's framing of working with messy real-world data and turning it into actionable insights, plus building interactive dashboards and web applications, sits at the same intersection where my work has been concentrated. At NYU's Novel AI Technologies I engineered \textbf{ETL pipelines} transforming \textbf{NoSQL} into \textbf{PostgreSQL} with concurrency control, then shipped \textbf{React, Node.js, REST APIs} dashboards for paying clients backed by row-level security and role-based access. At PwC I built \textbf{Power BI} dashboards over datasets joined across multiple relational databases for business stakeholders.

On statistical analysis and data mining: my Big Data graduate project built a \textbf{PySpark / HDFS} pipeline over millions of flight records, training Random Forest and Gradient Boosted Tree models with rigorous EDA, anomaly detection, and feature engineering. My coursework spans \textbf{Machine Learning}, \textbf{Deep Learning}, and \textbf{Big Data}.

On collaboration: I've partnered with product managers, business analysts, and marketing teams across PwC and NYU to translate ambiguous business questions into clear analytical work and shipped dashboards that real stakeholders use.

Stack-wise: \textbf{SQL}, \textbf{Python}, \textbf{Pandas}, \textbf{NumPy}, \textbf{scikit-learn}, \textbf{PyTorch}, plus \textbf{Power BI} and \textbf{QuickSight} for visualization. Comfortable with \textbf{React} and \textbf{Node.js} for building interactive web applications when the dashboard surface needs more than off-the-shelf BI tools.

What pulls me toward U-Michigan is the framing of "data for the common good." Philanthropy-driven engagement runs on actually understanding donor and program behavior in messy data, where careful, reproducible analytics compound into meaningful institutional outcomes.

I'd welcome the chance to discuss further. Thanks for considering my application.

\vspace{12pt}
Sincerely, \\
Ajay Venkatesh

\end{document}